
\section{Basic TFL Semantics}

TFL is \emph{truth-functional} as a result of its \emph{semantics}. Put simply,
we evaluate the \emph{truth} of a complex sentences as a \emph{function} of the
truth of their component parts. Ultimately this is a complex function of the
truth of all of its unique atomic sentences. We can define the notion with an
adaptation of the definition in \cite[][p. 72]{Mag25}

\begin{definition}[Truth-Functional]
  \label{def:truth-functional}
  A TFL connective is \emph{truth-functional} iff the truth value of the
  sentence comprising its scope (i.e. the sentence for which it is the main
  logical operator) is uniquely determined by the truth values of the
  sentence(s) it connects.
\end{definition}

Consider our informal definition of \enot\ephi.

\begin{argument}
\item[] A TFL sentence of the form \enot\ephi is true just in case \ephi is false. 
\end{argument}

We can see right away how this sort of principle works in practice. The rule is
simple. But actually working with TFL sentences within the system outlined by TFL
requires a substantial shift in how we think about these `sentences'. Indeed, even
this first definition hints at some important features of (classical) logic that
are different than those in ordinary English. We will discuss these in more detail
later, but one is that every TFL sentence is either true or false; if it is not
true, then it is false; if it is not false, then it is true. This rule (the Law of
the Excluded Middle) is exceptionless in TFL. We will formulate an axiom to that
effect later.

Recall that the process of symbolization for TFL loses a lot of the force or
affect that we often take for granted in ordinary language. This list of sentences
taken from \cite[][p. 73]{Mag25} all have the same truth-value and can be
symbolized with the same TFL sentence ($A \wedge B$), despite a variety of shaded
meanings.

\begin{enumerate}
  \item Dana is a logician and Dana is a nice person.
  \item Although Dana is a logician, Dana is a nice person.
  \item Dana is a logician despite being a nice person.
  \item Dana is a logician and Dana is a nice person.
  \item Dana is a nice person but also a logician. 
  \item Dana's being a logician notwithstanding, he is a nice person.
\end{enumerate}

The collapse of these diverse sentences to a single symbolization in TFL happens
precisely \textbf{because} its semantics is truth-functional. To see this, it will
help to look at the important terms one at a time: \emph{truth}, \emph{truth-value}
and \emph{function}.

So far we have used `truth' to mean something like a metalogical property that applies to
English sentences. On this informal usage, we can fall back on whatever folk definition we
like. This one from American Heritage Dictionary will do as well as any.

\begin{definition}[Informal Truth] \emph{noun}
  \begin{enumerate}
  \item Conformity to fact or actuality. \\ "Does this story have any truth?"
  \item Reality; actuality. \\ "In truth, he was not qualified for the job."
  \item The reality of a situation. \\ "The truth is, she respects your work."
  \end{enumerate}
  The American Heritage® Dictionary of the English Language, 5th Edition 
\end{definition}

We use this notion in understanding our symbolization keys. On our informal
treatment, the following key tells us to consider the truth of A and B to be a
matter of John and Mary's feelings about each other `in reality.'

\begin{argument}
\item[A:] John loves Mary.
\item[B:] Mary loves John.
\end{argument}

But in TFL, what the key really says is that $A$ has the same \emph{truth value}
that the English sentence `John loves Mary' does in reality. We are using `truth
value' here as a technical term.

\begin{definition}[Truth Value]
  A truth value is one of two values: \emph{True}) and \emph{False} (symbolized
  in the metalogic as `T' and `F', respectively. Every atomic sentence in TFL is
  assigned one or the other of these values and every complex sentence in TFL can
  be assigned one on the basis of its atomic parts.
\end{definition}

Recall from algebra that a \emph{function} takes one or more inputs
(\emph{arguments}) and produces an output (a \emph{result}). In TFL both the
arguments and the output of truth-functional connectives are truth values. Put
differently, the symbol for a TFL sentence is best understood \textbf{as} its
truth value. The meanings of the English sentences provided in the symbolization
key are just guides to assigning those values properly. Meaning plays no role in
TFL, only truth. This is an important point that takes some time to understand.

One consequnce of the truth-functional treatment of sentences and connectives is
that we can now fully define the semantics of a TFL connective on the basis of
how the truth of its complex sentence is determined by the truth of the sentences
it connects. We do this with the help of truth tables. Truth-tables exhaustively
spell out the consequences of assigning truth-values to the components of a
sentence. Note that characteristic truth tables use the metalogical variables
\ephi and \epsi because the rules can be applied to arbitrary TFL sentences, not
just to atomic ones.

\begin{definition}[Charateristic Truth Table: Negation]
  \label{def:tt-not}
  A TFL sentence of the form \enot\ephi is true just in case \ephi is false. The
  semantics of the negation operator `\enot' is defined by the following
  characteristic truth table: \\
  \begin{center}
    \begin{tabular}{c|c}
      \ephi & \enot\ephi\\
      \hline
      T & F\\
      F & T 
    \end{tabular}
  \end{center}
\end{definition}

For the negation operator, there is only one argument: the sentence negated; and
that sentence has only two possible truth values: $T$ and $F$.

\begin{definition}[Charateristic Truth Table: Conjunction]
  \label{def:tt-and}
  A TFL sentence of the form \ephi\eand\epsi is true just in case both \ephi and
  \epsi are true. The semantics of the connective `\eand' is defined by the
  following characteristic truth table: \\
  \begin{center}
    \begin{tabular}{c c|c}
      \ephi & \epsi & \ephi\eand\epsi\\
      \hline
      T & T & T\\
      T & F & F\\
      F & T & F\\
      F & F & F
    \end{tabular}
  \end{center}
\end{definition}

For the two-place connectives, there are two arguments: one for each of the
sentences connected; and each of those sentence has only two possible truth
values: $T$ and $F$. If we vary them independently for each sentence, we get
four total possibilities.\footnote{Note that the TTFF and TFTF ordering is
conventional for TFL truth tables. More generally, propositions to the left are
`chunked' and those on the right are `interleaved.' More on this later.}

\begin{definition}[Charateristic Truth Table: Disjunction]
  \label{def:tt-or}
  A TFL sentence of the form \ephi\eor\epsi is true just in case either \ephi is
  true or \epsi is true (or both). The semantics of the connective `\eor' is
  defined by the following characteristic truth table: \\
  \begin{center}
    \begin{tabular}{c c|c}
      \ephi & \epsi & \ephi\eor\epsi\\
      \hline
      T & T & T\\
      T & F & T\\
      F & T & T\\
      F & F & F
    \end{tabular}
  \end{center}
\end{definition}
\medskip
\begin{definition}[Charateristic Truth Table: Conditional]
  \label{def:tt-if}
  A TFL sentence of the form \ephi\eif\epsi is true just in case either \epsi is
  true or \ephi is false. The semantics of the connective `\eif' is defined by
  the following characteristic truth table: \\
  \begin{center}
    \begin{tabular}{c c|c}
      \ephi & \epsi & \ephi\eif\epsi\\
      \hline
      T & T & T\\
      T & F & F\\
      F & T & T\\
      F & F & T
    \end{tabular}
  \end{center}
\end{definition}
\medskip
\begin{definition}[Charateristic Truth Table: Biconditional]
  \label{def:tt-iff}
  A TFL sentence of the form \ephi\eiff\ephi is true just in case \ephi and \epsi
  are either both true or both false. The semantics of the connective `\eiff' is
  defined by the following characteristic truth table: \\
  \begin{center}
    \begin{tabular}{c c|c}
      \ephi & \epsi & \ephi\eiff\epsi\\
      \hline
      T & T & T\\
      T & F & F\\
      F & T & F\\
      F & F & T
    \end{tabular}
  \end{center}
\end{definition}

NOTE: Chapter 10 ends with a discussion of indicative versus subjunctive
conditionals. Add that somewhere (maybe later?)

Instead of depending on our knowledge of reality to assign truth-values to atomic
sentencs as we do in the informal assessment of truth in our symbolization tables,
we can simply assign values to each of the atomic sentence letters used in our
logic. Such an assignment is called a \emph{valuation}.

\begin{definition}[Valuation]
  \label{def:valuation}
  A \emph{valuation} is any assignment of truth values to the atomic sentence
  letters of TFL. On a valuation, each atomic sentence letter is assigned exaclty
  one truth value.
\end{definition}

Truth tables are \emph{exhaustive} in the sense that every possible valuation for a
given set of sentences can be represented as a row in the table, with $2^n$ such rows
for a TFL sentence with $n$ different atomic sentence letters. How this works in
practice is best demonstrated by example. Consider the TFL sentence $(A\wedge B)
\rightarrow A$. We can see how the truth of the whole depends on the truth of its
parts by working through each of the possible valuations `from the inside out'.

First, set up the table with all of the possible atomic truth values.

\begin{center}
  \begin{tabular}{c c|c c c c c}
    $A$ & $B$ & $(A$ & \eand & $B)$ & \eif & $A$ \\
    \hline
    T & T & & & & &\\
    T & F & & & & &\\
    F & T & & & & &\\
    F & F & & & & &
  \end{tabular}
\end{center}

Then add the atomic truth values to the appropriate columns on the right side of
the table.

\begin{center}
  \begin{tabular}{c c|c c c c c}
    $A$ & $B$ & $(A$ & \eand & $B)$ & \eif & $A$ \\
    \hline
    T & T & T & & T & & T\\
    T & F & T & & F & & T\\
    F & T & F & & T & & F\\
    F & F & F & & F & & F
  \end{tabular}
\end{center}

Now we can evaluate the truth of the sentences with connectives, starting with the
`innermost' ones (i.e. the ones that connect only atomic sentences). Here we start
with the \eand connective.

\begin{center}
  \begin{tabular}{c c|c c c c c}
    $A$ & $B$ & $(A$ & \eand & $B)$ & \eif & $A$ \\
    \hline
    T & T & T & \bf{T} & T & & T\\
    T & F & T & \bf{F} & F & & T\\
    F & T & F & \bf{F} & T & & F\\
    F & F & F & \bf{F} & F & & F
  \end{tabular}
\end{center}

We continue the process until we reach the main logical operator (i.e. the
outermost connective). The truth value assigned at this column is the truth value
for the entire sentence.

\begin{center}
  \begin{tabular}{c c|c c c c c}
    $A$ & $B$ & $(A$ & \eand & $B)$ & \eif & $A$   \\
        &     &      & \ephi &      & \eif & \epsi \\
    \hline
    T & T & T & T & T & \bf{T} & T\\
    T & F & T & F & F & \bf{T} & T\\
    F & T & F & F & T & \bf{T} & F\\
    F & F & F & F & F & \bf{T} & F
  \end{tabular}
\end{center}

Notice how we can use the metalogical sentence variables as substitutes for more
complex TFL sentences. Here they are really being used as \emph{logical} sentence
variables, representing not arbitrary sentences but these specific two:

\begin{argument}
  \item[\ephi] $A \wedge B$
  \item[\epsi] $A$
\end{argument}

Returning to our example from last week about the literacy of a certain professor,
we can see how this approach becomes complicated fast. Recall that the three
readings we found for our example were:

\begin{argument}
  \item[\ephi] $(A \eor \enot B) \eand C$
  \item[\epsi] $A \eor \enot (B \eand C)$
  \item[\echi] $A \eor (\enot B \eand C)$
\end{argument}

We can explore the different truth conditions for these sentences by expressing
each of them in the form of a truth table. Now that we have three atomic sentences
($A$, $B$, and $C$), we will need $2^3 = 8$ rows to express all of the possible
valuations relevant to this argument.
\footnote{Insert a discussion of r-possible and c-possible hereabouts.}
These valuations are straight-forward, if time-consuming, to lay out no matter how
many atomic sentences occur in our argument.

\begin{definition}[Complete Truth Table]
  The complete truth table for a TFL sentence \ephi is a truth table with each
  c-possible valuation; that is, for each possible combination of truth value
  assignments to all the unique atomic sentences found in \ephi.

  We can construct a complete truth table for a TFL sentence with $n$ unique atomic
  sentences using the following procedure:
  \begin{enumerate}
  \item Write each atomic sentence letter and the TFL sentence across the top,
    assuming a column for each of its connectives and atomic sentences.
  \item Insert $2^n$ rows.
  \item Populate the first column (the first $n$ columns are the unique atomic
    sentence letters) by adding `T' to the first half ($2^n/2$) of the rows and
    then false for the remaining half.
  \item For each subsequent column $i<=n$, fill in the portion of the column
    directly to the right of each uninterrupted series of Ts (or Fs) by filling
    in the first half with `T' and the second half with 'F'.
  \item If done correctly, the last column of atomic letters (column $n$) should
    alternate `T' then `F' $2^n/2$ times.
  \item Fill in the remaining columns according to the parse tree and connective
    semantics as above.
  \end{enumerate}

\end{definition}

Now that we know how to build a truth table for three atomic sentences, we can
build truth tables knowing that they are complete. First, set up the table with
all of the possible atomic truth values.

%$A$ & $B$ & $C$ & $A$  & \eor & \enot    & $(B$ & \eand & $C)$ \\
%$A$ & $B$ & $C$ & $(A$ & \eor & $(\enot$ & $B$  & \eand & $C$  \\
\begin{center}
  \begin{tabular}{c c c|c c c c c c}
    $A$ & $B$ & $C$ & $(A$ & \eor & \enot&$B)$ & \eand & $C$  \\
    \hline
    T & T & T & & & & & &\\
    T & T & F & & & & & &\\
    T & F & T & & & & & &\\
    T & F & F & & & & & &\\    
    F & T & T & & & & & &\\
    F & T & F & & & & & &\\
    F & F & T & & & & & &\\
    F & F & F & & & & & &
  \end{tabular}
\end{center}

Copy the truth values of the atomic sentences into the approriate columns on the
right side of the table.

\begin{center}
  \begin{tabular}{c c c|c c c c c c}
    $A$ & $B$ & $C$ & $(A$ & \eor & \enot&$B)$ & \eand & $C$  \\
    \hline
    T & T & T & T & & & T & & T\\
    T & T & F & T & & & T & & F\\
    T & F & T & T & & & F & & T\\
    T & F & F & T & & & F & & F\\    
    F & T & T & F & & & T & & T\\
    F & T & F & F & & & T & & F\\
    F & F & T & F & & & F & & T\\
    F & F & F & F & & & F & & F
  \end{tabular}
\end{center}

Apply the connective rules up the parse tree, starting with the atomic sentences.
\footnote{Starting here, we will use the convention from \cite{Smi20} where
multiple new steps are indicated in the same table by numbering along the bottom.}

\begin{center}
  \begin{tabular}{c c c|c c c c c c}
    $A$ & $B$ & $C$ & $(A$ & \eor & \enot&$B)$ & \eand & $C$  \\
    \hline
    T & T & T & T & T & F & T & \bf{T} & T\\
    T & T & F & T & T & F & T & \bf{F} & F\\ 
    T & F & T & T & T & T & F & \bf{T} & T\\
    T & F & F & T & T & T & F & \bf{F} & F\\    
    F & T & T & F & F & F & T & \bf{F} & T\\
    F & T & F & F & F & F & T & \bf{F} & F\\
    F & F & T & F & T & T & F & \bf{T} & T\\
    F & F & F & F & T & T & F & \bf{F} & F\\
    \hline
      &   &   &   & 2 & 1 &   & 3      &  \\
  \end{tabular}
\end{center}

The truth of the main sentence is provided for each valuation in the bolded column.
Similarly, the truth conditions for the other two sentences are expressed in the
bolded columns of the following two complete truth tables.

\begin{center}
  \begin{tabular}{c c c|c c c c c c}
    $A$ & $B$ & $C$ & $A$ & \eor & \enot & $(B$ & \eand & $C)$ \\
    \hline
    T & T & T & T & \bf{T} & F & T & T & T\\
    T & T & F & T & \bf{T} & T & T & F & F\\
    T & F & T & T & \bf{T} & T & F & F & T\\
    T & F & F & T & \bf{T} & T & F & F & F\\    
    F & T & T & F & \bf{F} & F & T & T & T\\
    F & T & F & F & \bf{T} & T & T & F & F\\
    F & F & T & F & \bf{T} & T & F & F & T\\
    F & F & F & F & \bf{T} & T & F & F & F\\
    \hline
      &   &   & 1 & \bf{6} & 5 & 2 & 4 & 3
  \end{tabular}
\end{center}

\begin{center}
  \begin{tabular}{c c c|c c c c c c}
    $A$ & $B$ & $C$ & $(A$ & \eor & $(\enot$ & $B$  & \eand & $C$  \\
    \hline
    T & T & T & T & \bf{T} & F & T & F & T\\ 
    T & T & F & T & \bf{T} & F & T & F & F\\
    T & F & T & T & \bf{T} & T & F & T & T\\
    T & F & F & T & \bf{T} & T & F & F & F\\    
    F & T & T & F & \bf{T} & F & T & F & T\\
    F & T & F & F & \bf{F} & F & T & F & F\\
    F & F & T & F & \bf{T} & T & F & T & T\\
    F & F & F & F & \bf{F} & T & F & F & F\\
    \hline
      &   &   & 1 & \bf{6} & 4 & 2 & 5 & 3
  \end{tabular}
\end{center}

If we compare the results in a single collapsed table, it is easy to see that these
TFL sentences have different truth conditions; that is, these sentences not true at
the same valuations. I've bolded the outliers to emphasize the differences.

\begin{center}
  \begin{tabular}{c c c|c c c}
    $A$ & $B$ & $C$ & \ephi & \epsi & \echi \\
    \hline
    T & T & T &     T  &     T   &     T \\
    T & T & F & \bf{F} &     T   &     T \\
    T & F & T &     T  &     T   &     T \\
    T & F & F & \bf{F} &     T   &     T \\    
    F & T & T &     F  &     F   & \bf{T}\\
    F & T & F &     F  & \bf{T}  &     F \\
    F & F & T &     T  &     T   &     T \\
    F & F & F &     F  & \bf{T}  &     F
  \end{tabular}
\end{center}

Next week, we will discuss the semantic implications of these differences in detail. And we
will look at some ways to simplify our process of constructing truth tables, which we will
continue into the following week.
